\section{NeRF: Implicit Neural Radiance Fields}
\label{sec:nerf}

Neural Radiance Fields occupy the foundational stratum of modern \nvs{}.
Their introduction by \citet{mildenhall2020nerf} demonstrated, for the first
time, that a compact neural parameterisation could represent complex real scenes
with view-dependent appearance at a quality that classical geometry-and-texture
pipelines could not approach.
The subsequent literature has pursued two independent but equally important
directions: \emph{accelerating} per-scene NeRF optimisation to make it
practical, and \emph{generalising} it beyond per-scene fitting to enable fast
inference across diverse, previously unseen scenes.
A third, deeply important direction addresses \emph{dynamic and deformable scenes},
which is a necessary prerequisite for world modeling.
We survey all three sub-families below, with Table~\ref{tab:nerf_methods}
providing a comprehensive cross-family comparison.

\subsection{Foundational NeRF Methods}

The original NeRF~\citep{mildenhall2020nerf} demonstrated that an MLP with
positional encoding could faithfully represent complex real scenes with
view-dependent appearance from a few dozen posed photographs.
Subsequent work systematically addressed its two principal limitations---slow
optimisation and poor generalisation---laying the groundwork for the modern
NeRF ecosystem.
Before examining these extensions, however, we first discuss improvements to
the core per-scene optimisation pipeline, as they established the quality and
speed baselines against which all later methods are measured.

\paragraph{Accelerating per-scene optimisation.}
The central bottleneck in vanilla NeRF is the MLP query: every rendered pixel
requires hundreds of network evaluations along a sampled ray, making training
convergence require hours even on high-end GPUs.
The insight that motivated the first generation of accelerated methods was that
the large MLP is architecturally overspecified for a fixed scene---much of its
capacity is consumed by storing spatial correlations that could be encoded
far more efficiently in an explicit spatial data structure.

Instant-NGP~\citep{muller2022instant} realised this insight most completely,
replacing the large MLP backbone with a compact multi-resolution
\emph{hash grid} that stores feature vectors at a trainable table of spatial
positions.
The hash encoding is both compact and highly parallelisable on GPU, enabling
training convergence in under a minute on a single GPU while matching or
exceeding prior quality.
TensoRF~\citep{chen2022tensorf} took a complementary algebraic approach,
decomposing the radiance field tensor into low-rank CP and VM components;
this tensor factorisation reduces both storage and memory bandwidth without
sacrificing the continuous differentiability of the implicit representation.
Plenoxels~\citep{fridovich2022plenoxels} eliminates the MLP entirely, using
a sparse 3-D voxel grid with trilinearly interpolated spherical harmonics;
the key observation is that view-dependent colour can be represented
without a neural network, and the resulting model trains in minutes by
avoiding costly backward passes through deep networks.
DirectVoxGO~\citep{sun2022dvgo} accelerates convergence via a two-stage
coarse-to-fine voxel grid that aggressively skips empty space during sampling,
directly reducing the number of expensive MLP queries to only those near
visible surfaces.

\paragraph{Improving visual quality and handling challenging scenes.}
Even as the acceleration methods were resolving the speed bottleneck,
a parallel strand of work addressed systematic quality failures in vanilla
NeRF: aliasing at different viewing distances, failure on unbounded scenes,
and poor handling of specular and metallic materials.

Mip-NeRF~\citep{barron2021mipnerf} identified aliasing as arising from the
point-sampling abstraction of the ray: when a pixel footprint covers a large
solid angle---as at large distances or with wide-angle lenses---querying the
radiance field at a single point is a poor approximation of the integral
over the pixel.
Mip-NeRF replaces the point-sampled ray with an
\emph{integrated positional encoding} over a conical frustum, substantially
reducing aliasing across multiple viewing distances and establishing the first
rigorous anti-aliasing framework for neural radiance fields.
Mip-NeRF~360~\citep{barron2022mipnerf360} extended this framework to
\emph{unbounded} outdoor scenes, which vanilla NeRF and even Mip-NeRF handle
poorly because the scene volume is infinite.
It resolves this via a non-linear scene parameterisation that contracts the
infinite background into a bounded representation, allowing the model to
allocate capacity appropriately between the foreground and the distant environment.
Zip-NeRF~\citep{barron2023zipnerf} later unified the hash grid of Instant-NGP
with the anti-aliasing of Mip-NeRF~360, establishing the current quality
frontier for per-scene NeRF by combining the speed of explicit encoding with
the rigour of multi-scale anti-aliasing.

NeRF-W~\citep{martinbrualla2021nerfw} addressed a different quality challenge:
reconstructing landmarks from uncontrolled internet images, where illumination
varies across photographs and transient objects (tourists, cars) are present in
some images but not others.
It introduces per-image latent appearance codes to absorb illumination variation,
together with a transient object model that down-weights inconsistent observations
during training.
Ref-NeRF~\citep{verbin2022refnerf} tackled the failure of directional encoding
to represent specular materials accurately, restructuring the directional
encoding into a physically-based reflection model that explicitly predicts
diffuse colour, specular tint, and surface roughness.
The key contribution is a \emph{reflectance decomposition} that separates the
predicted colour into diffuse and specular components, each with its own
directional dependence, dramatically improving the rendering of specular and
metallic surfaces.

\paragraph{Surface reconstruction.}
While volume rendering naturally produces accurate alpha composites, it does not
enforce a well-defined surface: the density function can be ``foggy'' near
object boundaries, leading to noisy geometry.
This matters greatly for applications beyond view synthesis---mesh extraction
for rendering in game engines, collision geometry for robotics, or structural
analysis in architecture---and motivated a dedicated line of work on
SDF-based neural surface reconstruction.

NeuS~\citep{wang2021neus} made the pivotal observation that volume rendering
and SDF reconstruction can be unified: by reformulating the volume rendering
integral using a logistic density derived from a Signed Distance Function,
it obtains watertight, high-quality surface reconstruction without any
explicit surface supervision.
VolSDF~\citep{yariv2021volsdf} independently derived a rigorous analytical
mapping from SDF values to volume density, providing a theoretically grounded
alternative and demonstrating that the logistic function used by NeuS is one
member of a broader family of valid density functions.
BakedSDF~\citep{yariv2023bakedsdf} takes the SDF-based surface quality one
step further by combining high-quality SDF reconstruction with real-time
rendering: after optimisation, the learned SDF is extracted into a mesh and
the view-dependent appearance is ``baked'' into a texture, enabling real-time
playback of neurally reconstructed scenes.

Together, these three sub-families---acceleration, quality, and
surface---resolved the core technical limitations of vanilla NeRF and set the
stage for the generalisation and dynamics work discussed next.

Table~\ref{tab:nerf_methods} provides a comprehensive comparison of NeRF variants
across all three sub-families: per-scene optimisation, generalizable,
and dynamic methods.

\begin{table*}[t]
  \centering
  \caption{Comprehensive NeRF method comparison.
    \textbf{Category}: PS=per-scene, GEN=generalizable, DYN=dynamic.
    \textbf{Train}/\textbf{Render}: speed relative to vanilla NeRF (Fast/Med/Slow).
    \textbf{Gen.}: generalizes across scenes without per-scene re-training.
    \textbf{Dyn.}: handles dynamic or deformable scenes.
    \textbf{Surf.}: produces explicit surface (mesh/SDF).}
  \label{tab:nerf_methods}
  \scriptsize
  \setlength{\tabcolsep}{3pt}
  \begin{tabular}{@{}llccccccp{4.2cm}@{}}
    \toprule
    \textbf{Method} & \textbf{Venue} & \textbf{Cat.}
      & \textbf{Train} & \textbf{Render}
      & \textbf{Gen.} & \textbf{Dyn.} & \textbf{Surf.}
      & \textbf{Key Contribution} \\
    \midrule
    \multicolumn{9}{@{}l}{\textit{Per-scene optimisation}} \\
    NeRF~\citep{mildenhall2020nerf}            & ECCV'20    & PS & Slow & Slow & \no & \no & \no & MLP + positional encoding \\
    NeRF-W~\citep{martinbrualla2021nerfw}       & CVPR'21    & PS & Slow & Slow & \no & \no & \no & In-the-wild appearance codes \\
    Mip-NeRF~\citep{barron2021mipnerf}         & ICCV'21    & PS & Slow & Slow & \no & \no & \no & Integrated positional encoding \\
    NeuS~\citep{wang2021neus}                  & NeurIPS'21 & PS & Slow & Slow & \no & \no & \yes & SDF-based volume rendering \\
    VolSDF~\citep{yariv2021volsdf}             & NeurIPS'21 & PS & Slow & Slow & \no & \no & \yes & Analytical SDF-to-density \\
    Plenoxels~\citep{fridovich2022plenoxels}   & CVPR'22    & PS & Fast & Fast & \no & \no & \no & No-MLP voxel grid + SH \\
    TensoRF~\citep{chen2022tensorf}            & ECCV'22    & PS & Med  & Fast & \no & \no & \no & CP/VM tensor decomposition \\
    DirectVoxGO~\citep{sun2022dvgo}            & CVPR'22    & PS & Fast & Fast & \no & \no & \no & Two-stage coarse-to-fine voxel \\
    Instant-NGP~\citep{muller2022instant}      & TOG'22     & PS & Fast & Fast & \no & \no & \no & Multi-resolution hash grid \\
    Mip-NeRF360~\citep{barron2022mipnerf360}   & CVPR'22    & PS & Slow & Slow & \no & \no & \no & Unbounded scene contraction \\
    Ref-NeRF~\citep{verbin2022refnerf}         & CVPR'22    & PS & Slow & Slow & \no & \no & \no & Physics-based reflection model \\
    Zip-NeRF~\citep{barron2023zipnerf}         & ICCV'23    & PS & Fast & Fast & \no & \no & \no & Hash grid + anti-aliasing \\
    BakedSDF~\citep{yariv2023bakedsdf}         & SIGGRAPH'23& PS & Med  & Fast & \no & \no & \yes & SDF mesh extraction + baking \\
    \midrule
    \multicolumn{9}{@{}l}{\textit{Generalizable methods}} \\
    pixelNeRF~\citep{yu2021pixelnerf}          & CVPR'21    & GEN& Fast & Slow & \yes& \no & \no & Pixel-aligned CNN features \\
    IBRNet~\citep{wang2021ibrnet}              & CVPR'21    & GEN& Fast & Slow & \yes& \no & \no & Cross-view transformer aggregation \\
    MVSNeRF~\citep{chen2021mvsnerf}            & ICCV'21    & GEN& Fast & Slow & \yes& \no & \no & Cost-volume geometry prior \\
    SRT~\citep{sajjadi2022srt}                 & CVPR'22    & GEN& Fast & Fast & \yes& \no & \no & Set-latent scene representation \\
    GNT~\citep{wang2022gnt}                    & arXiv'22   & GEN& Fast & Med  & \yes& \no & \no & Fully attention-based, no NeRF MLP \\
    NerfDiff~\citep{gu2023nerfdiff}            & ICML'23    & GEN& Fast & Slow & \yes& \no & \no & Diffusion prior for single-view \\
    ZeroRF~\citep{shi2023zerorf}               & arXiv'23   & GEN& Med  & Slow & \pt & \no & \no & Zero-pretrain sparse-view 360$^\circ$ \\
    FrameNeRF~\citep{yan2024framenerf}         & arXiv'24   & GEN& Fast & Slow & \yes& \no & \no & Video frames as pseudo views \\
    \midrule
    \multicolumn{9}{@{}l}{\textit{Dynamic \& deformable methods}} \\
    D-NeRF~\citep{pumarola2021dnerf}           & CVPR'21    & DYN& Slow & Slow & \no & \yes& \no & Canonical deformation network \\
    NSFF~\citep{li2021neural_scene_flow}       & CVPR'21    & DYN& Slow & Slow & \no & \yes& \no & Explicit scene flow field \\
    Nerfies~\citep{park2021nerfies}            & ICCV'21    & DYN& Slow & Slow & \no & \yes& \no & Per-frame latent deformation \\
    HyperNeRF~\citep{park2021hypernerf}        & TOG'21     & DYN& Slow & Slow & \no & \yes& \no & Higher-dimensional ambient space \\
    Non-Rigid NeRF~\citep{tretschk2021nonrigid}& ICCV'21    & DYN& Slow & Slow & \no & \yes& \no & Monocular ray bending \\
    DyNeRF~\citep{li2022neural3d}              & CVPR'22    & DYN& Slow & Slow & \no & \yes& \no & Time-conditioned multi-view NeRF \\
    TiNeuVox~\citep{fang2022tineuvox}          & SIGGRAPH'22& DYN& Med  & Med  & \no & \yes& \no & Voxel-accelerated dynamic NeRF \\
    K-Planes~\citep{fridovich2023kplanes}      & CVPR'23    & DYN& Fast & Fast & \no & \yes& \no & 4-D plane decomposition \\
    HexPlane~\citep{cao2023hexplane}           & CVPR'23    & DYN& Fast & Fast & \no & \yes& \no & Concurrent 4-D plane factorisation \\
    NeRFPlayer~\citep{song2023nerfplayer}      & TOG'23     & DYN& Med  & Med  & \no & \yes& \no & Static/deforming/new region decomp. \\
    Soup-of-Planes~\citep{du2023soupofplanes}  & ICCV'23    & DYN& Fast & Fast & \no & \yes& \no & Planar fitting for casual video \\
    Pseudo-GDV~\citep{zhao2024pgdvs}           & CVPR'24    & DYN& Med  & Slow & \no & \yes& \no & Static + dynamic inpainting \\
    \bottomrule
  \end{tabular}
\end{table*}

\subsection{Generalizable NeRF}

Per-scene NeRF achieves high quality but at the cost of lengthy per-scene
optimisation, which is fundamentally incompatible with applications requiring
rapid scene understanding---robotic manipulation, autonomous navigation, and
real-time AR all require scene understanding on the order of seconds, not hours.
\emph{Generalisable} methods address this bottleneck by learning a shared
prior across many training scenes, enabling fast inference on previously
unseen scenes without any per-scene optimisation.
The key design challenge is determining how to encode the reference views
into a scene representation that the NeRF decoder can query efficiently.

pixelNeRF~\citep{yu2021pixelnerf} established the foundational paradigm:
it conditions the NeRF MLP on pixel-aligned CNN features extracted from
reference views, so that at any query point along a ray, the model can
attend to the corresponding image-space features.
This pixel-aligned conditioning enables single-image reconstruction but
inherits the NeRF MLP's slow rendering speed.
IBRNet~\citep{wang2021ibrnet} improved upon this by using a transformer to
aggregate features from multiple reference views at query points along the ray,
learning to weight each reference view by its relevance to the current
query location and thus producing smoother, more consistent geometry.
MVSNeRF~\citep{chen2021mvsnerf} introduced explicit geometric reasoning
by building a cost volume from stereo correspondences, providing a
geometry-aware feature representation that substantially improves reconstruction
under large view baselines.

A particularly important direction is the elimination of the NeRF MLP
altogether.
GNT~\citep{wang2022gnt} dispenses with the volumetric rendering MLP,
replacing it with a two-stage transformer: the first stage encodes reference
views into a view-independent scene representation, and the second renders
the target view via cross-attention.
This architecture learns geometric reasoning as an emergent property of
attention rather than an explicit inductive bias, achieving better
generalisation on scenes that violate the assumptions of cost-volume methods.
SRT~\citep{sajjadi2022srt} pursues a similar philosophy with a
set-latent architecture, encoding an arbitrary number of views into a
fixed-size scene representation without pose assumptions, enabling
pose-free inference that anticipates the DUSt3R line of work.

More recent approaches have augmented generalisable NeRF with generative priors
to handle the under-constrained regime of very sparse input.
NerfDiff~\citep{gu2023nerfdiff} incorporates a diffusion model to synthesise
additional virtual views that supplement the sparse reference set, effectively
using the diffusion prior as a learned geometric hallucinator.
ZeroRF~\citep{shi2023zerorf} achieves competitive sparse-view 360$^\circ$
reconstruction via per-scene optimisation with a factorised feature grid and
no pretraining, demonstrating that careful regularisation can substitute for
generalisation in some sparse-input regimes.

\paragraph{Few-shot and test-time methods.}
The sparse-input regime has attracted sustained attention because it is the
most practically relevant: in deployment, reference images are scarce.
FrameNeRF~\citep{yan2024framenerf} exploits the temporal redundancy in
video by treating consecutive frames as dense pseudo-views, dramatically
increasing effective input density from monocular video without requiring
multi-camera hardware.
Recollection from Pensieve~\citep{zhou2024pensieve} takes a step further,
learning NVS from uncalibrated internet videos without any posed supervision,
treating pose estimation as an internal latent variable.
Test-Time Completion~\citep{bao2024tttcompletion} repurposes natural video
completion as a zero-shot NVS prior at test time, demonstrating that large
video models already encode sufficient spatial knowledge for competitive NVS
without additional training.
Mixed Neural Voxels~\citep{wang2023mixedvoxels} accelerates multi-view video
synthesis by decomposing scenes into static and dynamic voxel grids,
exploiting the temporal smoothness of most dynamic scenes.

Taken together, generalisable NeRF methods represent a conceptual shift from
the per-scene paradigm toward the \emph{amortised inference} view of neural
rendering: a single trained model encodes a universal geometric prior that
can be applied to any new scene at inference time.
This shift is a prerequisite for world modeling---a world model that required
hours of optimisation for each new scene would be useless in real-time decision-making
settings---and it directly motivates the feed-forward reconstruction models
discussed in Section~\ref{sec:generalisable}.

\subsection{Dynamic NeRF}

Modelling non-rigid dynamic scenes is a central challenge for world modeling,
as it requires the representation to capture not just spatial geometry but
also the temporal evolution of scene structure.
The fundamental difficulty is that a dynamic scene lacks a fixed canonical
frame: the geometry changes continuously, and a naive extension of static
NeRF to incorporate time as an additional input dimension often fails to
model large deformations faithfully.

The first generation of dynamic NeRF methods addressed this through the
\emph{canonical deformation} paradigm.
D-NeRF~\citep{pumarola2021dnerf} introduces a deformation network that maps
each spacetime query point $(x,t)$ to a corresponding location in a canonical
static space, where a standard NeRF computes appearance; at test time, novel
views at any timestamp are rendered by composing the deformation and canonical
fields.
Nerfies~\citep{park2021nerfies} refines this idea for casual monocular video
of human subjects, where the long-range correlations of body motion make
pure deformation network predictions unreliable; it instead uses per-frame
latent deformation codes with regularisation to prevent physically implausible
deformations.
HyperNeRF~\citep{park2021hypernerf} generalises further to scenes with
topological changes---mouth opening, paper tearing---which require changes
in surface connectivity that smooth deformations cannot represent.
It embeds the canonical space in a higher-dimensional ambient space, allowing
the ``slice'' seen at each timestep to represent scenes of changing topology.

Alongside the deformation paradigm, Neural Scene Flow Fields~(NSFF)~\citep{li2021neural_scene_flow}
proposed an orthogonal approach: explicitly modelling the 3-D velocity field
of the scene, enabling supervision from monocular depth estimation and optical
flow and providing an intuitive physical interpretation of motion.
Non-Rigid NeRF~\citep{tretschk2021nonrigid} reconstructs and synthesises
novel views of dynamic scenes from monocular video with explicit ray bending,
demonstrating that accurate reconstruction from a single moving camera is
achievable even for complex non-rigid motion.
DyNeRF~\citep{li2022neural3d} extends NeRF to multi-view video by conditioning
on time, providing a cleaner experimental setting with strong geometric
grounding from simultaneous cameras.

The second generation of dynamic NeRF methods focused on efficiency, motivated
by the observation that per-scene dynamic NeRF optimisation is even slower
than its static counterpart.
TiNeuVox~\citep{fang2022tineuvox} accelerates dynamic NeRF by replacing the
deformation-field MLP with a time-conditioned voxel feature grid, reducing
training time from days to hours.
The key architectural innovation came with K-Planes~\citep{fridovich2023kplanes}
and HexPlane~\citep{cao2023hexplane}, which independently proposed the same
elegant solution: decompose 4-D space-time into six
feature planes---three spatial~(XY, XZ, YZ) and three spatial-temporal
(XT, YT, ZT)---each encoded as a learned 2-D grid, making the architecture
highly memory-efficient and interpretable.
Their simultaneous publication demonstrated the generality of the plane-based
4-D decomposition idea and established it as the dominant acceleration paradigm
for dynamic NeRF.
NeRFPlayer~\citep{song2023nerfplayer} took a complementary approach to efficiency
by separately encoding static, deforming, and newly-appearing scene regions,
allowing the model to allocate capacity according to local motion statistics.

More recent work has pushed toward casual, monocular dynamic video synthesis.
Soup-of-Planes~\citep{du2023soupofplanes} enables fast view synthesis of
casual monocular videos by fitting a lightweight set of planar scene
representations; this simple parameterisation is surprisingly effective
for the forward-facing camera motions common in casual video.
Pseudo-Generalised Dynamic View~\citep{zhao2024pgdvs} synthesises dynamic
novel views from monocular video by combining static reconstruction with
per-frame dynamic inpainting, separating the estimation problem into
tractable sub-problems.
Decoupled Dynamic Video~\citep{wang2024decoupleddynamic} disentangles
foreground dynamics and background statics, assigning different model
capacities to each and enabling cleaner decomposition of casual monocular videos.

The trajectory of dynamic NeRF research mirrors that of static NeRF: from
proof-of-concept demonstrations requiring dense camera rigs to fast,
practical methods operating on casual monocular video.
This trajectory is directly connected to the world-modeling agenda: a world
model must reason about dynamic scenes from single-camera observations,
precisely the regime where the most recent dynamic NeRF methods operate.
The dynamic NeRF methods are included in the comprehensive
Table~\ref{tab:nerf_methods} above.

% ─────────────────────────────────────────────────────────────────────────────
